\section{Conclusões}

\subsection{Resultados alcançados}

Destaca-se que o sistema desenvolvido atingiu os objetivos propostos, demonstrando eficácia na gestão do fluxo de peças entre os tapetes. A implementação das abordagens centralizadas e distribuídas permitiu uma análise comparativa detalhada, evidenciando as vantagens e desvantagens de cada método em diferentes cenários operacionais. Os resultados obtidos indicam que a abordagem distribuída, especialmente em regime assíncrono, proporcionou maior flexibilidade e adaptabilidade às variações no fluxo de peças, apesar das complexidades adicionais introduzidas. A análise dos atrasos de comunicação revelou a importância de mecanismos robustos para garantir a integridade dos dados e a eficiência do sistema. Acrescenta-se ainda que as simulações realizadas forneceram insights valiosos sobre o comportamento do sistema sob diferentes condições, contribuindo para a validação das escolhas de design e implementação adotadas ao longo do projeto.

\subsection{Dificuldades encontradas e soluções adotadas}

Durante o desenvolvimento do sistema, foram enfrentadas diversas dificuldades técnicas e conceituais. A implementação da comunicação entre módulos na abordagem distribuída apresentou desafios significativos, especialmente na gestão dos atrasos de comunicação. Para mitigar esses problemas, foram adotados protocolos de comunicação robustos e mecanismos de buffer para garantir a integridade dos dados transmitidos. Além disso, a coordenação entre os módulos em regime síncrono exigiu a implementação de algoritmos de sincronização eficientes, que foram desenvolvidos e testados exaustivamente para assegurar o correto funcionamento do sistema. A complexidade adicional introduzida pela abordagem assíncrona também demandou soluções inovadoras para garantir a consistência do estado do sistema, o que foi alcançado através da utilização de técnicas de controle distribuído e monitoramento contínuo do fluxo de peças. Essas soluções permitiram superar os desafios encontrados, resultando em um sistema robusto e eficiente, capaz de operar de forma eficaz em diferentes cenários e condições operacionais.